%% notes-tex/025-additive-baselines/sections/4-ladder.tex

\section{The ladder under both splits\secstatus{todo}}
\label{sec:ladder}

%% SOURCE: experiments/025-solid-growth/scripts/additive_baselines_025_panels.py ->
%% tables/t2-heldout.tex, from results/additive_baselines_025.csv, the summary json,
%% and 010-kuzmin-tmi/results/additive_baseline_gene_interaction.csv.
\begin{table}[htbp]
\centering
\footnotesize
\caption{Test Pearson for every model under both arms, beside the 010 numbers. B0
through B4 are closed-form fits on a fixed split and carry no spread; B5 is the mean
and standard deviation over three seeds. The three transformer rows stand on
different footings, named in the row label: the 010 checkpoints were scored on the
010 test split of 37{,}673 records; GH 1598 and GH 1640 are validation maxima over
36 logged epochs, and neither has a test score.}
\label{tab:heldout}
%% GENERATED by experiments/025-solid-growth/scripts/additive_baselines_025_panels.py
%% SOURCE: the result files named in that script's docstring. Do not edit by hand.
\begin{tabular}{lrrr}
\toprule
 & \multicolumn{2}{c}{Arm R, random over records} & Arm Q, query-pair disjoint \\
Model & 010 build & 025 build & 025 build \\
\midrule
B0 train mean & 0.000 & 0.000 & 0.000 \\
B4 query pair only & 0.390 & 0.390 & 0.000 \\
B3 screen mean & 0.390 & 0.390 & 0.142 \\
B1 additive per-gene ridge & 0.400 & 0.400 & 0.185 \\
B2 additive plus pair ridge & 0.406 & 0.406 & 0.180 \\
B5 embedding MLP, 3 seeds & $0.426 \pm 0.001$ & $0.432 \pm 0.005$ & $0.141 \pm 0.007$ \\
CGT, three 010 checkpoints, test & 0.438 to 0.455 & & \\
CGT GH 1598, val max, epoch 14 of 36 & & 0.446 & \\
CGT GH 1640, val max, epoch 7 of 36 & & & 0.199 \\
\bottomrule
\end{tabular}

\end{table}

%% SOURCE: experiments/025-solid-growth/scripts/additive_baselines_025_panels.py ->
%% additive_baselines_025_fig1_ladders.svg. Panels a, b: the test rows of
%% results/additive_baselines_025.csv plus the transformer references in the summary
%% json. Panel c: results/additive_baselines_025_arm_val_history.csv. Panel d:
%% 010-kuzmin-tmi/results/query_pair_disjoint_cv.csv against the arm Q test rows.
\tcfig{figures/additive_baselines_025_fig1_ladders.pdf}{The ladder under both
splits. \textbf{a}, arm R, random over records: the six baselines on the 37{,}673
test records, the three 010 checkpoints on the 010 test split, and the 025
replication GH 1598 as a validation maximum, hatched. \textbf{b}, arm Q, query-pair
disjoint, on the 37{,}791 test records of 46 held-out screens, with GH 1640 hatched.
In both panels only B5 carries an error bar, the standard deviation over its three
seeds; every other bar is one deterministic fit. \textbf{c}, validation Pearson per
epoch for the two arms of the 010 configuration, each against the additive ridge fit
on its own validation part; open circles mark each run's best epoch, the value the
hatched bars carry. \textbf{d}, how far the disjoint null moves with the choice of
held-out screens: five disjoint folds on the 010 build per model, mean and standard
deviation with each fold as a point, and the arm Q single split as a diamond. B0 and
B4 are zero on every fold and are omitted.}{fig:ladders}

\notesec{Arm R is the 010 result}
Every closed-form model agrees with its 010 value to $3 \times 10^{-13}$, which
follows from the record-level label and split identity of Section~\ref{sec:data}
rather than testing it. B5 differs by 0.006 against a three-seed standard deviation
of 0.005; the 025 fit ran on CPU where 010's ran on GPU, and the cause is untraced.
GH 1598 reaches a validation maximum of 0.446 at epoch 14, within 0.001 of the lower
end of the 0.447 to 0.462 band the three 010 checkpoints recorded on validation.

\notesec{On the random split the ladder is mostly screen identity}
B4 knows only which query pair a triple carries and reaches 0.390, 97 percent of the
additive ridge and 86 percent of the best checkpoint. The transformer's margin over
B1 is real, $+0.038$ to $+0.055$ with paired-bootstrap intervals that exclude zero
in the 010 analysis, and it is small.

\notesec{Holding out screens collapses the ladder and inverts it}
On arm Q the additive ridge falls from 0.400 to 0.185 and B4 to zero by
construction. B5 falls from 0.432 to 0.141, below the ridge: the capacity that
helped on unseen records of seen screens hurts on unseen screens. B3 falls back from
a screen mean to a gene mean and lands at 0.142.

\notesec{One disjoint split is a small sample of screens}
Panel d is the reason arm Q's numbers are read as one draw rather than as the
disjoint null. Across five disjoint folds on the 010 build the ridge spans 0.088 to
0.151, $0.127 \pm 0.033$, and arm Q's 0.185 lies above every fold. The same holds
for B2, B3 and B5, so arm Q is a comparatively easy draw of 46 held-out screens. A
transformer number on arm Q is compared to the arm Q nulls, on the same 46 screens,
and not to the fold mean.
